\begin{abstract}
Traditional goal-oriented dialogue agents often require costly retraining when
new objectives or constraints arise. Recent multi-objective approaches improve
this limitation by parameterizing goals, but they typically rely on explicit
numeric preference vectors. This interface is poorly aligned with practical
settings, where users express macro-goals in natural language and expect the
agent to balance strategic flexibility with non-negotiable constraints. The
problem becomes harder under intent drift: as the dialogue context evolves, a
static policy may continue following an outdated short-term strategy even when
the user's macro-goal calls for adaptation. We propose \textbf{ALMOND}, a
hierarchical framework for learning adaptive dialogue policies from
natural-language macro-goals. ALMOND combines a low-level policy that executes
communication skills with a regret-based mechanism that prioritizes
under-mastered skills, enabling complex behaviors to be bootstrapped from
simpler ones. A high-level policy monitors intent shifts and selects
short-term skills using strategic hints distilled during training, so that
local behavior remains aligned with the overarching linguistic goal. We
evaluate ALMOND under constrained buyer-side negotiation with
persona-conditioned seller drift. Across benchmark settings, ALMOND improves
LLM-as-Judge goal completion by 4--7 percentage points, reduces adaptation
delay by 2--3 turns, and lowers constraint violations compared with
static-weight and ablation baselines. Phase-1 skill training also improves
$r$-fair, $r$-gain, and $r$-deal by 3--6 percentage points while reducing
average turns.
\end{abstract}
